\begin{pastexam}[label={px:ncu-111-dsa-05}]{中央資工所 111 學年度 複選題第 5 題}
（複選；每題 5 分，全對才給分，答錯倒扣 1 分）

Which of the following statements are false?
\begin{enumerate}[label=(\Alph*)]
\item The function \textbf{F1} below is executable (i.e., terminable) for all positive integer x.
\begin{minted}[fontsize=\small]{text}
Int F1(int x)
{
     If x is even then
        return x/2;
     else
        return F1(F1(3x+2));
}
\end{minted}
\item When input data are 21 and 12, the output of function \textbf{F2} below is 4.
\begin{minted}[fontsize=\small]{text}
Int F2(int x, int y)
{
     If y=0 then
        return x;
     else
        return F2(y, x mod y);
}
\end{minted}
\item Consider the function \textbf{F3} below. \textbf{F3}(5) is 5.
\begin{minted}[fontsize=\small]{text}
Int F3(int x)
{
     Int p, q;

     If x <= 2 then
        return x;
     else
        p=F3(x-2);
        q=F3(x-3);
        return p+q;
}
\end{minted}
\item None of above.
\end{enumerate}
\end{pastexam}

\begin{pastsoln}
答案是 (A)、(B)、(C)、(D)，四個都是錯誤的敘述。

三段程式是虛擬碼（\texttt{If \ldots{} then}、\texttt{3x} 表示 $3 \times x$），照字面的意思追蹤
（資結書第~6.4~節）。
\begin{itemize}
\item (A) 錯誤。 $x$ 是奇數時 $3x$ 是奇數，$3x + 2$ 也是奇數，所以內層的 \texttt{F1(3x+2)} 又走到
\texttt{else}，再呼叫 \texttt{F1(3(3x+2)+2)}，參數仍是奇數而且越來越大，永遠到不了 \texttt{x/2} 那一支。
例如 $x = 1$：\texttt{F1(1)} 呼叫 \texttt{F1(5)}，\texttt{F1(5)} 呼叫 \texttt{F1(17)}，\texttt{F1(17)}
呼叫 \texttt{F1(53)}……不會結束（實際執行會用完呼叫堆疊）。 只有偶數 $x$ 會結束，所以「對所有正整數
都會結束」是錯的。
\item (B) 錯誤。 \texttt{F2} 是輾轉相除法（資結書虛擬碼~6.9）：
\[
\mathtt{F2}(21, 12) = \mathtt{F2}(12, 9) = \mathtt{F2}(9, 3) = \mathtt{F2}(3, 0) = 3,
\]
輸出是 $\gcd(21, 12) = 3$，不是 $4$。
\item (C) 錯誤。 $\mathtt{F3}(0) = 0$、$\mathtt{F3}(1) = 1$、$\mathtt{F3}(2) = 2$，
\[
\mathtt{F3}(3) = \mathtt{F3}(1) + \mathtt{F3}(0) = 1, \qquad
\mathtt{F3}(5) = \mathtt{F3}(3) + \mathtt{F3}(2) = 1 + 2 = 3 \ne 5.
\]
（$\mathtt{F3}(5)$ 只用到 $x \ge 0$ 的呼叫，不會碰到負數。）
\item (D) 錯誤。 「None of above」在「哪些是錯的」這一題裡表示「上面沒有錯的」，但 (A)、(B)、(C) 都是
錯的，所以這句本身也錯。
\end{itemize}
\end{pastsoln}
